% ============================================================
% UNIDAD 6: CABLES AERONÁUTICOS, MAZOS Y EWIS
% ============================================================
\clearpage
\section*{UNIDAD 6: CABLES AERONÁUTICOS}
\subsection*{Tipos de cable, selección, mazos, blindajes, masas y EWIS}

\begin{recuadro}
\textbf{Objetivo de la unidad:} conocer los tipos de cables usados en aeronaves, seleccionar la
sección adecuada, armar e instalar mazos según las prácticas aceptadas, y aplicar correctamente
los blindajes y las masas.
\end{recuadro}

% ------------------------------------------------------------
\section*{1. POR QUÉ NO CUALQUIER CABLE}

En una aeronave los cables soportan vibración, temperaturas extremas, fluidos (combustible,
aceite, fluido hidráulico) y, en caso de incendio, no deben propagar la llama ni producir humo
tóxico. Por eso se usan cables con normas específicas. \textbf{Nunca se usan cables con aislación
de PVC} (los de uso domiciliario o automotor): emiten gases tóxicos y corrosivos al quemarse.

El conjunto de cables, conectores, terminales, soportes y protecciones de una aeronave se denomina
\textbf{EWIS} (\textit{Electrical Wiring Interconnection System}). Después de accidentes como el
del TWA 800 (1996) y el Swissair 111 (1998), en los que el cableado fue protagonista, las
autoridades reforzaron los requisitos de diseño, inspección y mantenimiento del cableado.

% ------------------------------------------------------------
\section*{2. TIPOS DE CABLES}

\begin{center}
\small
\begin{tabular}{|p{0.28\linewidth}|p{0.34\linewidth}|p{0.26\linewidth}|}
\hline
\textbf{Norma} & \textbf{Características} & \textbf{Uso} \\ \hline
SAE AS22759 (ex MIL-W-22759) /16, /32, /34 & Conductor de cobre estañado, aislación de ETFE
(Tefzel) o ETFE reticulado; 600 V; hasta 150 $^\circ$C. & El más usado en aviónica de aviación
general \\ \hline
AS22759 con conductor plateado o niquelado & Mayor temperatura (200 $^\circ$C o más). & Zonas
calientes, motor \\ \hline
MIL-W-81381 (poliimida, ``Kapton'') & Muy delgado y liviano, pero se fisura y puede producir
\textbf{arcos que se propagan} (\textit{arc tracking}). & Aeronaves antiguas; hoy se evita y se
inspecciona con cuidado \\ \hline
M27500 & Cable multiconductor con o sin \textbf{blindaje} y cubierta, armado con conductores
AS22759. & Señales, audio, buses de datos \\ \hline
Coaxiales M17 (RG-400, RG-142) & 50 $\Omega$, PTFE. & Antenas (Unidad de microondas) \\ \hline
Cables de par trenzado especiales & Impedancia controlada (por ejemplo, para Ethernet o buses de
datos). & Redes de aviónica \\ \hline
\end{tabular}
\end{center}

\paragraph*{Cómo se lee una designación.} Por ejemplo, \texttt{M22759/16-22-9}: norma 22759,
hoja de especificación 16 (ETFE, cobre estañado), sección \textbf{22 AWG}, color \textbf{9}
(blanco). Un M27500 lleva en su designación la cantidad de conductores, su tipo, el blindaje y la
cubierta.

% ------------------------------------------------------------
\section*{3. SELECCIÓN DE LA SECCIÓN (AWG)}

La sección del cable se elige por dos criterios: que \textbf{no se caliente} (corriente admisible,
según si va solo o en mazo y la altitud) y que la \textbf{caída de tensión} no supere los límites.
Las tablas y gráficos de la circular \textbf{FAA AC 43.13-1B, capítulo 11}, son la referencia
habitual.

\begin{center}
\small
\begin{tabular}{|c|c|c|}
\hline
\textbf{AWG} & \textbf{Resistencia aprox. ($\Omega$ cada 1000 pies)} & \textbf{Uso típico en
aviónica} \\ \hline
24 & 25,7 & Señales, datos (cuando el fabricante lo permite) \\ \hline
22 & 16,1 & Señales, datos, audio, discretas \\ \hline
20 & 10,1 & Alimentación de equipos de bajo consumo \\ \hline
18 & 6,4 & Alimentación de equipos \\ \hline
16 & 4,0 & Alimentación de equipos de mayor consumo \\ \hline
14 & 2,5 & Barras secundarias, luces \\ \hline
12 & 1,6 & Barras, cargas grandes \\ \hline
10 & 1,0 & Alimentaciones principales \\ \hline
\end{tabular}
\end{center}

\paragraph*{Caída de tensión máxima recomendada} (AC 43.13-1B): en sistemas de 14 V, 0,5 V en
servicio continuo y 1 V intermitente; en sistemas de 28 V, 1 V continuo y 2 V intermitente.

\begin{recuadro}
\textbf{Ejemplo.} Una pantalla consume 3 A a 14 V y está a 4 m (13 pies) del disyuntor. El circuito
tiene ida y vuelta (o masa cercana): longitud efectiva $\approx$ 26 pies con cable 22 AWG:
$R = 26 \times 16{,}1 / 1000 = 0{,}42\ \Omega$; $\Delta V = 3 \times 0{,}42 = 1{,}26$ V. Supera los
0,5 V. Con 18 AWG: $R = 26 \times 6{,}4/1000 = 0{,}17\ \Omega$; $\Delta V = 0{,}5$ V. Con 18 AWG
queda justo en el límite; se elige 18 AWG (o 16 si se quiere margen), siempre respetando lo que
indica el manual de instalación del equipo.
\end{recuadro}

% ------------------------------------------------------------
\section*{4. IDENTIFICACIÓN DE LOS CABLES}

Cada cable se identifica con un código (número de circuito, segmento, sección) impreso sobre la
aislación (por láser o por estampado en caliente) o en manguitos termocontraíbles. Según la AC
43.13-1B, la identificación se coloca a intervalos regulares (no más de 15 pulgadas) y cerca de
cada extremo. El código debe coincidir con el \textbf{plano de cableado}: es lo que permite seguir
un circuito años después.

% ------------------------------------------------------------
\section*{5. MAZOS DE CABLES E INSTALACIÓN}

\begin{itemize}
  \item \textbf{Armado en tablero:} el mazo se arma sobre un tablero con el dibujo a escala y clavos
        o soportes que marcan el recorrido, las ramificaciones y las longitudes.
  \item \textbf{Atado:} con \textbf{hilo de atar encerado o de fibra} (nudos de amarre) o con
        precintos aprobados y apretados con la herramienta adecuada (sin cortar dejando filos).
  \item \textbf{Soporte:} \textbf{abrazaderas acolchadas} (tipo MS21919) cada cierta distancia,
        sin apretar ni aplastar los cables.
  \item \textbf{Pasos por estructura:} siempre con \textbf{pasacables} (\textit{grommets}) o
        protectores de borde; nunca un cable apoyado sobre un borde metálico.
  \item \textbf{Radio de curvatura:} no menor a 10 veces el diámetro exterior del cable o mazo mayor
        (con excepciones en los extremos, según la AC 43.13-1B); en coaxiales, el que indique el
        fabricante.
  \item \textbf{Separación:} los cables deben quedar separados de las líneas de combustible,
        aceite, oxígeno y fluidos, y de los comandos de vuelo; los cables de señales sensibles
        (audio, datos, magnetómetro) separados de los de potencia.
  \item \textbf{Holgura de servicio:} dejar suficiente longitud para poder desconectar y reparar
        un conector; evitar tensiones y roces con piezas móviles.
\end{itemize}

% ------------------------------------------------------------
\section*{6. BLINDAJES Y MASAS}

\begin{itemize}
  \item El \textbf{blindaje} (malla alrededor de los conductores) protege las señales de la
        interferencia y evita que los cables irradien.
  \item Se conecta a masa \textbf{donde lo indique el fabricante}: en un extremo (típico en audio,
        para evitar lazos de masa que producen zumbido) o en ambos (típico en buses de datos y RF).
  \item La conexión se hace con un \textbf{manguito de soldadura} (\textit{solder sleeve}) que suelda
        un cable de drenaje a la malla, o terminando la malla en el \textit{backshell} del conector.
        El ``rabo'' (\textit{pigtail}) de masa debe ser lo más corto posible.
  \item \textbf{Masas de alimentación:} cada equipo con su retorno a un punto de masa o a una barra
        de masas, con terminales bien ajustados sobre superficies limpias. Una masa deficiente
        genera fallas intermitentes, ruido y lecturas erróneas de motor.
  \item \textbf{Bonding} (unión eléctrica de la estructura): se mide en miliohms entre el equipo o
        su bandeja y la estructura.
\end{itemize}

% ------------------------------------------------------------
\section*{7. ACTIVIDADES DE LA UNIDAD}
\begin{enumerate}
  \item Interpretar las designaciones: \texttt{M22759/16-20-2} y \texttt{M22759/32-22-9}.
  \item Calcular la sección necesaria para un equipo de 5 A en 28 V a 6 m del disyuntor.
  \item Identificar en muestras de cables cuáles son aptos para uso aeronáutico y cuáles no.
  \item Armar en un tablero un pequeño mazo de tres ramas con atado de hilo y abrazaderas.
  \item Investigar el rol del cableado en el accidente del Swissair 111 y qué cambió después.
\end{enumerate}
